\section{From VNF to CNF: what the platform must provide}

Section~2 introduced the 5G core. This section does the work the original
framing was missing: it derives, from the concrete needs of that core,
\emph{exactly} which capabilities the NFVI platform must offer --- and then
maps each capability onto a specific Kubernetes mechanism. Every design choice
in Sections~4 and~5 traces back to a row of Table~\ref{tab:requirements}.

\subsection{Why containers (CNFs) and not VMs (VNFs)}

The first NFV deployments ran network functions as virtual machines (VNFs) on
OpenStack. That worked, but each VM carries a full guest OS: boot times of
tens of seconds, hundreds of megabytes to gigabytes of RAM before the function
even starts, and multi-gigabyte images. For a 5G core that must scale quickly
with demand, this is too heavy. A \textbf{CNF} re-packages the same function as
containers: it starts in seconds, shares the host kernel, ships as a compact
image, and --- most importantly --- delegates its lifecycle (restart, scale,
heal) to Kubernetes controllers instead of bespoke automation. Open5GS ships
each NF as such a container, which is why the whole core fits on a 16~GB host.

\subsection{Requirements derived from the 5G core}

Table~\ref{tab:requirements} is the heart of the report's argument. The
left column is a concrete fact about the 5G core; the middle column is the
platform capability it forces; the right column is how we provide it.

\begin{table}[H]
\centering
\small
\begin{tabular}{p{0.34\textwidth} p{0.30\textwidth} p{0.28\textwidth}}
\toprule
\textbf{Need of the 5G core} & \textbf{Required platform capability} & \textbf{Provided by} \\
\midrule
UPF forwards user data over dedicated N3/N6 interfaces, separate from signalling
  & A pod must have \emph{multiple} NICs on \emph{isolated} L2 networks
  & Multus + macvlan + Whereabouts (Sec.~4.2, 5.3) \\
Control NFs (NRF, AUSF\ldots) are stateless and scale with signalling load
  & On-demand horizontal scaling
  & HorizontalPodAutoscaler + metrics-server (Sec.~5.4) \\
NFs are independent functions with their own lifecycle
  & Declarative deploy / upgrade / heal per function
  & Kubernetes Deployments + Helm (Sec.~5.4) \\
NFs discover each other over SBI (HTTP/2)
  & Stable virtual IPs and L4 load balancing
  & Kubernetes Service (ClusterIP) \\
An operator must see the core's health and the UPF's load
  & Cluster-wide metrics and dashboards
  & Prometheus + Grafana (Sec.~5.6) \\
Subscriber records (IMSI, keys) must persist
  & A stateful datastore
  & MongoDB (UDR backend) \\
Management, signalling, and user traffic must not mix
  & Three isolated network planes
  & Three-plane design (Sec.~4.2) \\
\bottomrule
\end{tabular}
\caption{Each platform capability exists to satisfy a specific 5G-core need.}
\label{tab:requirements}
\end{table}

\subsection{Mapping ETSI / 5G onto Kubernetes}

Table~\ref{tab:mapping} fixes the vocabulary used throughout: how each ETSI NFV
block, and each concrete 5G artefact, is realised in this testbed.

\begin{table}[H]
\centering
\begin{tabular}{lll}
\toprule
\textbf{ETSI NFV block} & \textbf{5G artefact} & \textbf{This testbed (k3s)} \\
\midrule
NFVI            & Compute + network for NFs   & VirtualBox VMs node0/1/2 \\
VIM             & Resource manager            & k3s API server + scheduler \\
VNF / CNF       & UPF, AMF, SMF, NRF, \ldots   & Pods in namespace \texttt{open5gs} \\
VNFM            & NF lifecycle manager        & Helm \\
NFVO            & Service orchestrator        & Ansible + kubectl (manual) \\
Virtual network & N2/N3/N4/N6 + SBI           & Multus/macvlan + Flannel \\
Storage         & Subscriber DB               & MongoDB (emptyDir, testbed scope) \\
\bottomrule
\end{tabular}
\caption{Mapping of ETSI NFV blocks and 5G artefacts onto the k3s testbed.}
\label{tab:mapping}
\end{table}

\subsection{The multi-NIC mechanism in detail}

Because multi-NIC networking is the capability that makes the 5G user plane
possible, it is worth stating precisely. A standard pod gets one interface
(\texttt{eth0}) from the primary CNI (Flannel). \textbf{Multus} is a
\emph{meta}-CNI: it does not move packets itself but lets a pod attach
additional interfaces by delegating to other plugins. Here Multus delegates the
extra interfaces to \textbf{macvlan}, which creates virtual L2 interfaces bound
to a physical NIC (\texttt{eth2}). Addresses on these networks are handed out by
\textbf{Whereabouts}, a cluster-wide IPAM plugin that guarantees
non-overlapping allocation without a DHCP server. Each extra network is
described by a \texttt{NetworkAttachmentDefinition} (NAD) and requested by a pod
through the \texttt{k8s.v1.cni.cncf.io/networks} annotation. This is exactly the
machinery the UPF uses to obtain its N3 and N6 interfaces.
